\lecture{15}{2026-10-19}{Algebraic Generating Functions}

Suppose \(F(z)\) is an algebraic generating function satisfying \(P(z, F(z)) = 0\), where
\[
    P(z, y) = P_d(z) y^d + P_{d-1}(z) y^{d-1} + \cdots + P_0(z) \in \mathbb{Q}[z][y]
\]
with \(P_d(z) \neq 0\).
Using transfer theorems, Puiseux expansions, and the Newton--Puiseux method, we want to determine the asymptotic behaviour of the coefficients of \(F(z)\).
This raises two main questions:
\begin{itemize}
    \item Where are the singularities of the branches of \(P(z, y) = 0\)?
    \item How do we find the singularities and local expansions of our generating function \(F(z)\) (rather than one of its conjugates)?
\end{itemize}

\section*{Singularities of algebraic branches}

To address the first question, fix \(z = a \in \mathbb{C}\).
When are all solutions \(y(z)\) of \(P(z, y) = 0\) analytic around \(z = a\)?

By the implicit function theorem, if \(P(a, y)\) has \(d\) \emph{distinct} roots in \(y\), then all \(d\) branches are analytic around \(z = a\) \cite[Section 2.3, p.~46]{Melczer2021}.
This condition can fail in two ways:
\begin{enumerate}[label=\Alph*.]
    \item If \(P(a, y)\) has degree less than \(d\) in \(y\), which happens when \(P_d(a) = 0\). In this case, we say that a branch escapes to infinity.
    \item If \(P(a, y)\) has a multiple root in \(y\), which happens when \(P(a, y)\) and \(P_y(a, y)\) share a common root. In this case, we say that two or more branches collide.
\end{enumerate}

\begin{definition}[Resultant, {\cite[Definition 2.14, Problem 2.14]{Melczer2021}}]
    A tool to detect common roots between two polynomials \(f(y)\) and \(g(y)\) is the \emph{resultant} \(\operatorname{Res}(f, g)\), defined by
    \[
        \operatorname{Res}(f, g) = c_f^{\deg g} c_g^{\deg f} \prod_{\alpha, \beta} (\alpha - \beta),
    \]
    where \(c_f\) and \(c_g\) are the leading coefficients of \(f\) and \(g\), and the product is taken over all roots \(\alpha\) of \(f\) and \(\beta\) of \(g\) (counted with multiplicity).
    Equivalently, \(\operatorname{Res}(f, g)\) can be computed as the determinant of the \emph{Sylvester matrix} of \(f\) and \(g\), which is a \((\deg f + \deg g) \times (\deg f + \deg g)\) matrix whose entries are the coefficients of \(f\) and \(g\).
\end{definition}

In particular, \(f\) and \(g\) have a common root if and only if \(\operatorname{Res}(f, g) = 0\).

\begin{lemma}[{\cite[Definition 2.15, Lemma 2.2]{Melczer2021}}]
    If a branch of \(P(z, y) = 0\) has a singularity at \(z = a\), then
    \[
        P_d(a) = 0 \quad \text{or} \quad \operatorname{Res}_y(P, P_y)(a) = 0.
    \]
    In particular, when \(P \in \mathbb{Q}[z, y]\), only algebraic numbers can appear as singularities of algebraic generating functions.
\end{lemma}

\begin{corollary}[{\cite[Corollary 2.1]{Melczer2021}}]
    Let \(F(z) = \sum_{n \geq 0} f_n z^n\) be an algebraic series analytic at the origin, with dominant singularities \(\omega_0\beta, \ldots, \omega_{m-1}\beta\) where \(\beta \in \mathbb{R}_{>0}\) and \(|\omega_j| = 1\) for all \(j\). Then
    \[
        f_n = \frac{\beta^{-n} n^s}{\Gamma(s + 1)} \sum_{j = 0}^{m - 1} c_j \omega_j^{-n} + O\left(\beta^{-n} n^{t}\right),
    \]
    where \(s \in \mathbb{Q} \setminus \{-1, -2, \ldots\}\), \(t < s\), and \(\beta, \omega_j, c_j\) are algebraic numbers.
\end{corollary}

\begin{example}[{\cite[Example 2.14]{Melczer2021}}]
    The algebraic generating function
    \[
        F(z) = (1-4z)^{-1/2}
    \]
    has coefficient sequence
    \[
        [z^n] F(z) = \binom{2n}{n} \sim \frac{4^n}{\sqrt{\pi n}}.
    \]
    Thus, the series \(A(z) = \sum_{n \geq 0} \binom{2n}{n}^2 z^n\) is transcendental, since
    \[
        [z^n] A(z) = \binom{2n}{n}^2 \sim \frac{16^n}{\pi n}.
    \]
\end{example}